%! TEX root = kriteria_tapak_temu2.tex
% the main TeX file which is intended to compile, :VimtexReload after adjustment
   % this file should be included in the main file
% :h vimtex-tex-root
\documentclass[../main.tex]{subfiles}
%ensure the class of docs

\begin{document}

% ==========================================
% TOPIK 13A: UNSUR-UNSUR PERENCANAAN (Sub 1 & 2)
% ==========================================

\section{definisi tapak}%
\label{sec:definisi tapak}

\begin{frame}[mytitle=standard, t, mycolor=digiPH_gray, light]
	\frametitle{Unsur-Unsur Perencanaan Tapak}
	\framesubtitle{Bagian 1: Akses dan Organisasi Ruang}

	Unsur tapak sangat luas dan saling berhubungan.

	\vspace{0.3cm}
	\textbf{Akses dan Sirkulasi}
	\begin{itemize}
		\item \textit{Main entrance} dan \textit{side entrance}
		\item Sirkulasi dan area parkir
	\end{itemize}

	\vspace{0.3cm}
	\textbf{Organisasi Ruang}
	\begin{itemize}
		\item Tata guna lahan dan \textit{zoning}
		\item Pola ruang luar
		\item Bukaan dan tutupan ruang luar
	\end{itemize}
\end{frame}

% ==========================================
% TOPIK 13B: UNSUR-UNSUR PERENCANAAN (Sub 3 & 4)
% ==========================================
\begin{frame}[mytitle=standard, t, mycolor=digiPH_cream, light]
	\frametitle{Unsur-Unsur Perencanaan Tapak}
	\framesubtitle{Bagian 2: Massa Bangunan dan Utilitas}

	\textbf{Massa dan Bangunan}
	\begin{itemize}
		\item Bentuk massa dan bentuk bangunan
		\item Orientasi dan komposisi massa
		\item Bukaan dan tutupan massa
	\end{itemize}

	\vspace{0.3cm}
	\textbf{Kenyamanan dan Utilitas}
	\begin{itemize}
		\item Penghawaan dan penerangan
		\item Jaringan komunikasi
	\end{itemize}

	\vspace{0.3cm}
	\textit{Semua unsur tersebut perlu dipertimbangkan sebagai satu kesatuan, bukan sebagai komponen yang terpisah.}
\end{frame}

% ==========================================
% TOPIK 14: FAKTOR ALAM (ANGIN)
% ==========================================
\begin{frame}[mytitle=standard, c, mycolor=digiPH_ocean, dark]
	\frametitle{Faktor Alam: Angin}

	\begin{description}
		\item[Angin laminer] Aliran angin berlapis-lapis dengan arah dan kecepatan yang relatif tidak berubah antar lapisan.
		\item[Angin terpisah] Terjadi ketika perubahan topografi menyebabkan perubahan pola aliran dan pemisahan lapisan angin.
		\item[Angin turbulen] Perubahan pola aliran menghasilkan turbulensi, terutama akibat perubahan topografi (bukit dan lembah).
	\end{description}

	\vspace{0.5cm}
	\textbf{Implikasi dalam desain:} \\
	Kontur + arah angin + posisi bangunan = memengaruhi kualitas aliran udara pada tapak.
\end{frame}

% ==========================================
% TOPIK 15A: FAKTOR ALAM (TANAH)
% ==========================================
\begin{frame}[mytitle=standard, t, mycolor=digiPH_white, light]
	\frametitle{Faktor Alam: Tanah dan Topografi}
	\framesubtitle{A. Analisis Tanah}

	Hal-hal yang perlu dianalisis terkait kondisi tanah pada tapak:
	\begin{itemize}
		\item Jenis dan kondisi tanah.
		\item Perubahan jenis tanah dalam tapak.
		\item Tingkat keasaman atau kebasaan.
		\item Ketebalan lapisan humus.
		\item Kemampuan tanah dalam menyerap air.
		\item Kemampuan tanah dalam mencegah erosi.
	\end{itemize}
\end{frame}

% ==========================================
% TOPIK 15B: FAKTOR ALAM (TOPOGRAFI)
% ==========================================
\begin{frame}[mytitle=standard, t, mycolor=digiPH_gray, light]
	\frametitle{Faktor Alam: Tanah dan Topografi}
	\framesubtitle{B. Analisis Topografi}

	Hal-hal yang perlu diperhatikan terkait topografi tapak:
	\begin{itemize}
		\item Kecuraman atau kedataran lereng.
		\item Keseragaman permukaan tanah.
		\item Hubungan dengan permukaan di sekitar tapak.
		\item Elemen alam yang tidak dapat diubah.
		\item Risiko terjadinya erosi.
		\item Orientasi lereng.
	\end{itemize}

	\vspace{0.3cm}
	\textbf{Kesimpulan:} Tanah dan topografi memengaruhi letak bangunan, sirkulasi, ruang luar, dan pengolahan permukaan tapak.
\end{frame}

% ==========================================
% TOPIK 16: KONTUR & CUT AND FILL
% ==========================================
\begin{frame}[mytitle=standard, c, mycolor=digiPH_leaf, dark]
	\frametitle{Kontur dan Konsep Cut and Fill}

	Kontur menggambarkan kondisi dan relief permukaan tapak.
	\begin{itemize}
		\item \textbf{CUT (KUPASAN):} Proses penggalian atau pemotongan tanah.
		\item \textbf{FILL (URUGAN):} Proses penimbunan atau penambahan tanah.
		\item \textbf{CUT AND FILL:} Kombinasi pemotongan dan penambahan tanah untuk memperoleh ketinggian permukaan yang sesuai.
	\end{itemize}

	\vspace{0.3cm}
	\textbf{Pertanyaan desain:} ``Apakah seluruh tapak harus dibuat datar?'' \\
	Tidak selalu. Tujuannya bukan sekadar meratakan tanah, tetapi menghasilkan kondisi tapak yang sesuai fungsi dan lingkungannya.
\end{frame}

% ==========================================
% TOPIK 17A: INTEGRASI (ALUR BERPIKIR)
% ==========================================
\begin{frame}[mytitle=center, t, mycolor=digiPH_white, light]
	\frametitle{Integrasi: Dari Analisis Menuju Perancangan}
	\framesubtitle{Alur Berpikir Perancangan Tapak}

	\begin{itemize}
		\item \textbf{BACA TAPAK:} Identifikasi kondisi alam, fisik, dan sosial.
		\item \textbf{TEMUKAN POTENSI \& KENDALA:} Apa yang dapat dimanfaatkan? Apa yang harus dihindari?
		\item \textbf{ANALISIS KAIDAH TAPAK:} Orientasi — Aksesibilitas — Vegetasi — View — Eksisting.
		\item \textbf{ANALISIS FAKTOR ALAM:} Angin — Tanah — Topografi — Kontur.
		\item \textbf{SUSUN PROGRAM \& ZONING:} Publik — Semipublik — Privat — Servis.
		\item \textbf{INTEGRASIKAN:} Bangunan + sirkulasi + ruang luar + utilitas + lingkungan sekitar.
	\end{itemize}
\end{frame}

% ==========================================
% TOPIK 17B: INTEGRASI (HASIL & REFLEKSI)
% ==========================================
\begin{frame}[mytitle=standard, c, mycolor=digiPH_cream, light]
	\frametitle{Integrasi: Dari Analisis Menuju Perancangan}
	\framesubtitle{Hasil Akhir dan Refleksi}

	\textbf{HASIL AKHIR:} \\
	Tapak yang fungsional, nyaman, aman, sehat, estetis, dan terpadu.

	\vspace{0.8cm}
	\begin{alertblock}{Refleksi Akhir}
		\centering
		\textit{“Perancangan tapak yang baik bukan dimulai dari menggambar bangunan, tetapi dari memahami tempat.”}
	\end{alertblock}
\end{frame}

\section{Kriteria Pemilihan Tapak}%
\label{sec:Kriteria Pemilihan Tapak}

\begin{frame}[mytitle=center,t,mycolor=digiPH_gray,light]


\end{frame}

% ---------------------------------------------------
% FRAME 1: Judul Utama & Pokok Bahasan
% ---------------------------------------------------
\begin{frame}[mytitle=center, c, mycolor=digiPH_ocean, dark]
	\frametitle{Kriteria Pemilihan Tapak}
	\framesubtitle{Dasar Analisis dalam Perancangan Arsitektur dan Lingkungan Binaan}

	\textbf{Pokok bahasan:}
	\begin{itemize}
		\item Lokasi dan luas tapak
		\item Bentuk tapak
		\item Topografi
		\item Aksesibilitas
		\item Kondisi tapak
		\item Ketampakan
		\item Tapak ideal
		\item Penilaian dan matriks pemilihan tapak
	\end{itemize}

	\vspace{0.5cm}
	\textbf{Gagasan utama:}\\
	Pemilihan tapak bukan sekadar memilih lokasi yang tersedia, tetapi menentukan tingkat kesesuaian lahan terhadap fungsi, bangunan, akses, kondisi fisik, dan tujuan proyek.
\end{frame}

% ---------------------------------------------------
% FRAME 2: Mengapa Pemilihan Tapak Penting?
% ---------------------------------------------------
\begin{frame}[mytitle=standard, t, mycolor=digiPH_gray, light]
	\frametitle{Mengapa Pemilihan Tapak Penting?}

	Pemilihan tapak merupakan keputusan awal yang memengaruhi berbagai keputusan perancangan berikutnya.

	\vspace{0.3cm}
	\textbf{Tapak menentukan:}
	\begin{itemize}
		\item Bagaimana bangunan ditempatkan.
		\item Bagaimana pengguna masuk dan keluar kawasan.
		\item Bagaimana bangunan merespons topografi dan kondisi tanah.
		\item Bagaimana vegetasi dipertahankan atau dimanfaatkan.
		\item Seberapa efektif lahan dapat dikembangkan.
		\item Potensi keberhasilan fungsi bangunan dari sisi lokasi.
	\end{itemize}

	\vspace{0.3cm}
	\textbf{Prinsip analitis:}
	\begin{itemize}
		\item Tapak yang baik bukan hanya tapak yang ``bagus secara visual'', tetapi tapak yang paling sesuai dengan kebutuhan proyek.
		\item Materi sumber menekankan bahwa lokasi yang dipilih perlu sesuai dengan hasil survei ekonomi dan tujuan proyek.
	\end{itemize}
\end{frame}

% ---------------------------------------------------
% FRAME 3: Faktor Utama
% ---------------------------------------------------
\begin{frame}[mytitle=standard, c, mycolor=digiPH_white, light]
	\frametitle{Faktor Utama dalam Pemilihan Tapak}

	Enam faktor utama yang perlu dianalisis secara bersama:

	\vspace{0.3cm}
	\begin{table}[]
		\centering
		\small
		\begin{tabular}{@{}lp{10cm}@{}}
			\toprule
			\textbf{Faktor}         & \textbf{Pertanyaan Analisis}                                     \\ \midrule
			\textbf{Lokasi \& luas} & Apakah lokasi strategis dan luasnya mencukupi?                   \\
			\textbf{Bentuk tapak}   & Apakah bentuknya rasional dan mudah dikembangkan?                \\
			\textbf{Topografi}      & Bagaimana kemiringan dan kontur memengaruhi bangunan?            \\
			\textbf{Aksesibilitas}  & Bagaimana hubungan tapak dengan jaringan jalan dan transportasi? \\
			\textbf{Kondisi tapak}  & Bagaimana vegetasi dan kondisi tanahnya?                         \\
			\textbf{Ketampakan}     & Bagaimana kualitas visual dan keberadaan tapak dalam lingkungan? \\ \bottomrule
		\end{tabular}
	\end{table}

	\vspace{0.3cm}
	\textit{Keenam faktor tersebut merupakan kerangka dasar yang disebutkan langsung dalam materi sumber.}
\end{frame}

% ---------------------------------------------------
% FRAME 4: Lokasi dan Luas
% ---------------------------------------------------
\begin{frame}[mytitle=standard, t, mycolor=digiPH_cream, light]
	\frametitle{Lokasi dan Luas Tapak}

	\textbf{Lokasi tapak}
	\begin{itemize}
		\item Harus mempertimbangkan fungsi dan tujuan proyek.
		\item Lokasi yang ideal bergantung pada karakter kegiatan.
		\item Untuk bangunan komersial, lokasi yang ramai dapat menjadi pertimbangan penting.
	\end{itemize}

	\textbf{Luas tapak}
	\begin{itemize}
		\item Tidak cukup hanya untuk \textit{footprint} bangunan.
		\item Harus mempertimbangkan: bangunan utama, fasilitas pendukung, ruang luar, sirkulasi, serta kemungkinan pengembangan.
	\end{itemize}

	\vspace{0.3cm}
	\textbf{Implikasi perancangan:}\\
	Tapak yang terlalu kecil dapat membatasi fleksibilitas, sedangkan tapak yang memadai memberikan ruang bagi pengembangan kawasan. Materi menyatakan bahwa luas tapak harus lebih besar daripada luasan bangunan.
\end{frame}

% ---------------------------------------------------
% FRAME 5: Bentuk
% ---------------------------------------------------
\begin{frame}[mytitle=standard, t, mycolor=digiPH_white, light]
	\frametitle{Bentuk Tapak}

	\textbf{Prinsip dasar:} Bentuk tapak harus rasional dan menguntungkan proses perencanaan.

	\vspace{0.3cm}
	\textbf{Hal yang dianalisis:}
	\begin{itemize}
		\item Keteraturan bentuk.
		\item Efisiensi pemanfaatan lahan.
		\item Hubungan antara batas tapak dan massa bangunan.
		\item Kemungkinan pembagian zona \& pembentukan sirkulasi.
		\item Hubungan bentuk tapak dengan ruang luar.
	\end{itemize}

	\vspace{0.3cm}
	\textbf{Pertanyaan analitis mahasiswa:}\\
	Apakah bentuk tapak mendukung konsep bangunan, atau justru memaksa desain untuk beradaptasi secara tidak efisien?

	\vspace{0.2cm}
	\footnotesize{\textit{Rujukan: Materi mengacu pada FIA UK (2018) dalam pembahasan karakter bentuk tapak.}}
\end{frame}

% ---------------------------------------------------
% FRAME 6: Topografi dan Kontur
% ---------------------------------------------------
\begin{frame}[mytitle=standard, t, mycolor=digiPH_leaf, dark]
	\frametitle{Topografi dan Kontur}

	\textbf{Topografi} = kondisi permukaan tanah yang memengaruhi proses perancangan bangunan.

	\vspace{0.3cm}
	\textbf{Dampak topografi terhadap desain:}
	\begin{itemize}
		\item Penempatan massa bangunan.
		\item Ketinggian bangunan terhadap lingkungan.
		\item Arah sirkulasi dan pengolahan ruang luar.
		\item Pekerjaan tanah dan hubungan bangunan dengan kontur alami.
	\end{itemize}

	\textbf{Strategi desain:}
	\begin{itemize}
		\item Menyesuaikan bangunan dengan kontur.
		\item Memanfaatkan perbedaan elevasi (termasuk \textit{cut and fill} bila perlu).
	\end{itemize}

	\vspace{0.3cm}
	\textbf{Prinsip:} Semakin baik desain merespons kondisi kontur, semakin kecil kecenderungan tapak dipaksa mengikuti bentuk bangunan.
\end{frame}

% ---------------------------------------------------
% FRAME 7: Sintesis Bentuk & Topografi
% ---------------------------------------------------
\begin{frame}[mytitle=standard, c, mycolor=digiPH_white, light]
	\frametitle{Hubungan Bentuk dan Topografi dengan Perancangan}

	\textbf{Bentuk tapak + topografi $\rightarrow$ konfigurasi bangunan}

	\vspace{0.3cm}
	\textbf{Contoh hubungan analitis:}
	\begin{itemize}
		\item \textbf{Tapak relatif teratur + topografi sederhana}\\
		      $\rightarrow$ penataan massa lebih fleksibel $\rightarrow$ pembagian zona lebih mudah $\rightarrow$ sirkulasi lebih sederhana.
		\item \textbf{Tapak tidak teratur + topografi kompleks}\\
		      $\rightarrow$ diperlukan strategi penempatan massa $\rightarrow$ kemungkinan lebih banyak pekerjaan tanah $\rightarrow$ desain harus lebih responsif.
	\end{itemize}

	\vspace{0.3cm}
	\textbf{Pertanyaan studio:}\\
	Apakah lebih tepat mengubah tapak untuk mengikuti bangunan, atau mengubah bangunan untuk mengikuti tapak? (Keduanya harus dipandang sebagai proses yang saling memengaruhi).
\end{frame}

% ---------------------------------------------------
% FRAME 8: Aksesibilitas
% ---------------------------------------------------
\begin{frame}[mytitle=center, c, mycolor=digiPH_navyblue, dark]
	\frametitle{Aksesibilitas Tapak}

	Aksesibilitas berkaitan dengan kemampuan tapak untuk menerima, menghubungkan, dan mengatur pergerakan pengguna.

	\vspace{0.3cm}
	\textbf{Aspek utama:}
	\begin{itemize}
		\item Jaringan jalan di sekitar tapak.
		\item Jalan masuk dan keluar.
		\item Hubungan dengan tapak lain.
		\item Akses transportasi umum dan kendaraan pribadi.
		\item Kesesuaian akses dengan fungsi bangunan.
	\end{itemize}

	\vspace{0.3cm}
	\textbf{Implikasi:}\\
	Lokasi yang strategis tetapi sulit dicapai belum tentu menjadi tapak yang paling sesuai. Jaringan jalan harus memberikan manfaat maksimal bagi fungsi yang direncanakan.
\end{frame}

% ---------------------------------------------------
% FRAME 9: Vegetasi (Menggunakan imageplus jika ada logo/ikon)
% ---------------------------------------------------
\begin{frame}[mytitle=imageplus, t, mycolor=digiPH_gray, light]
	\frametitle{Kondisi Tapak: Vegetasi}

	Vegetasi bukan sekadar elemen estetika, tetapi bagian dari sistem tapak.

	\vspace{0.3cm}
	\textbf{Fungsi vegetasi menurut materi:}
	\begin{itemize}
		\item Mengikat tanah pada permukaan tanah yang curam.
		\item Membantu menyerap air.
		\item Menjadi peneduh dan menghalau kebisingan.
	\end{itemize}

	\textbf{Implikasi desain:}
	\begin{itemize}
		\item Vegetasi eksisting perlu diidentifikasi sejak tahap analisis.
		\item Vegetasi yang memiliki fungsi ekologis atau fungsional dapat dipertahankan.
		\item Penempatan bangunan mempertimbangkan keberadaan vegetasi penting.
	\end{itemize}

	\vspace{0.2cm}
	\textbf{Prinsip:} Vegetasi eksisting adalah aset tapak, bukan semata-mata hambatan.
\end{frame}

% ---------------------------------------------------
% FRAME 10: Kondisi Tanah
% ---------------------------------------------------
\begin{frame}[mytitle=standard, t, mycolor=digiPH_white, light]
	\frametitle{Kondisi Tanah}

	Kondisi tanah memiliki hubungan langsung dengan perancangan struktur. Tanah yang keras dianggap lebih ideal karena mempermudah rancangan struktur.

	\vspace{0.3cm}
	\textbf{Hal yang perlu diperhatikan:}
	\begin{itemize}
		\item Karakter tanah eksisting.
		\item Kesesuaian tanah dengan kebutuhan bangunan.
		\item Implikasi terhadap sistem struktur dan pekerjaan tapak.
	\end{itemize}

	\vspace{0.3cm}
	\textbf{Hubungan sebab-akibat:}\\
	Kondisi tanah $\rightarrow$ memengaruhi keputusan struktur $\rightarrow$ memengaruhi konstruksi $\rightarrow$ memengaruhi biaya dan kelayakan pengembangan.

	\vspace{0.2cm}
	\footnotesize{\textit{Catatan: Analisis tapak awal tidak seharusnya menggantikan investigasi tanah/geoteknik lanjutan.}}
\end{frame}

% ---------------------------------------------------
% FRAME 11: Ketampakan
% ---------------------------------------------------
\begin{frame}[mytitle=standard, t, mycolor=digiPH_cream, light]
	\frametitle{Ketampakan dan Karakter Visual Tapak}

	Ketampakan berkaitan dengan bagaimana tapak terlihat, hadir, dan berhubungan dengan lingkungan sekitarnya.

	\vspace{0.3cm}
	\textbf{Hal yang dapat diamati:}
	\begin{itemize}
		\item Kualitas visual tapak.
		\item Hubungan tapak dengan elemen alami (garis pohon, topografi).
		\item Hubungan visual dengan lingkungan sekitar.
		\item Karakter ruang dan lanskap.
		\item Potensi menjadi bagian dari identitas bangunan/kawasan.
	\end{itemize}

	\vspace{0.3cm}
	\textbf{Prinsip:}\\
	Tapak yang baik memungkinkan bangunan membentuk hubungan yang harmonis dengan karakter alami dan lingkungan sekitarnya.
\end{frame}

% ---------------------------------------------------
% FRAME 12: Tapak Ideal
% ---------------------------------------------------
\begin{frame}[mytitle=standard, c, mycolor=digiPH_darkblue, dark]
	\frametitle{Karakteristik Tapak yang Ideal}

	Berdasarkan materi, tapak ideal memiliki beberapa karakter utama:

	\begin{enumerate}
		\item \textbf{Lokasi sesuai tujuan proyek}\\
		      Lokasi mendukung fungsi dan kelayakan proyek.
		\item \textbf{Luas mencukupi}\\
		      Mampu menampung bangunan, fasilitas pendukung, serta pengembangan.
		\item \textbf{Mendukung pembangunan yang harmonis}\\
		      Bangunan dapat beradaptasi dengan karakter alami tapak.
		\item \textbf{Memperhatikan penerimaan publik}\\
		      Pemilihan lokasi mempertimbangkan kesesuaian sosial terhadap proyek.
	\end{enumerate}

	\vspace{0.3cm}
	\footnotesize{\textit{Peringatan: Memulai proyek tanpa mempertanyakan kesesuaian lokasi dan penerimaan publik merupakan kesalahan dalam perencanaan.}}
\end{frame}

% ---------------------------------------------------
% FRAME 13: Kriteria Penilaian Ekonomi
% ---------------------------------------------------
\begin{frame}[mytitle=standard, t, mycolor=digiPH_gray, light]
	\frametitle{Kriteria Penilaian Tapak}

	Pemilihan tapak tidak hanya dilakukan melalui pengamatan fisik, tetapi juga dapat menggunakan kriteria penilaian valuasi.

	\vspace{0.3cm}
	\textbf{Enam kriteria penilaian dalam materi:}
	\begin{itemize}
		\item Perbandingan penjualan (\textit{Sales Comparison})
		\item Alokasi (\textit{Allocation})
		\item Ekstraksi (\textit{Extraction})
		\item Pembagian pembangunan (\textit{Subdivision Development})
		\item Nilai sisa tanah (\textit{Land Residual})
		\item Kapitalisasi sewa dasar (\textit{Ground Rent Capitalization})
	\end{itemize}

	\vspace{0.3cm}
	\textbf{Perspektif mahasiswa:}\\
	Kriteria tersebut menunjukkan bahwa kualitas tapak dapat dipertimbangkan dari sudut pandang ekonomi dan nilai lahan, bukan hanya kualitas fisiknya.
\end{frame}

% ---------------------------------------------------
% FRAME 14: Matriks Bagian 1 (Konsep Tahapan)
% ---------------------------------------------------
\begin{frame}[mytitle=standard, t, mycolor=digiPH_white, light]
	\frametitle{Matriks Pemilihan Tapak}
	\framesubtitle{Langkah Sistematis}

	Ketika tersedia beberapa alternatif lokasi, keputusan perlu dibuat secara sistematis menggunakan matriks pemilihan.

	\vspace{0.3cm}
	\textbf{Tahapan:}
	\begin{enumerate}
		\item \textbf{Menentukan alternatif tapak:} (A, B, C, dst.)
		\item \textbf{Menentukan kriteria:} Lokasi, luas, bentuk, topografi, aksesibilitas, kondisi tapak, ketampakan.
		\item \textbf{Menetapkan ukuran/nilai penilaian:} Setiap kriteria diberikan nilai berdasarkan tingkat kualitas lahannya.
		\item \textbf{Memberikan skor} pada setiap alternatif.
		\item \textbf{Membandingkan hasil}.
		\item \textbf{Menentukan tapak} dengan tingkat kesesuaian terbaik.
	\end{enumerate}
\end{frame}

% ---------------------------------------------------
% FRAME 15: Matriks Bagian 2 (Tabel Ilustrasi)
% ---------------------------------------------------
\begin{frame}[mytitle=standard, c, mycolor=digiPH_white, light]
	\frametitle{Matriks Pemilihan Tapak}
	\framesubtitle{Contoh Struktur Matriks}

	Hubungan antara kriteria dan alternatif dilakukan melalui tabel pemilihan, dengan tingkat kualitas lahan dinyatakan dalam angka.

	\vspace{0.3cm}
	\begin{table}[]
		\centering
		\small
		\begin{tabular}{@{}lcccc@{}}
			\toprule
			\textbf{Kriteria} & \textbf{Bobot} & \textbf{Tapak A} & \textbf{Tapak B} & \textbf{Tapak C} \\ \midrule
			Lokasi            & 20\%           & 4                & 5                & 3                \\
			Luas              & 15\%           & 3                & 5                & 4                \\
			Bentuk            & 15\%           & 4                & 3                & 5                \\
			Topografi         & 15\%           & 5                & 3                & 4                \\
			Aksesibilitas     & 20\%           & 4                & 5                & 3                \\
			Kondisi tapak     & 15\%           & 4                & 4                & 3                \\ \bottomrule
		\end{tabular}
	\end{table}

	\vspace{0.2cm}
	\footnotesize{\textit{Catatan: Angka pada contoh di atas hanya untuk ilustrasi metode penjumlahan matriks.}}
\end{frame}

% ---------------------------------------------------
% FRAME 16: Sintesis (Bagian 1 dari Sintesis & Tugas)
% ---------------------------------------------------
\begin{frame}[mytitle=center, c, mycolor=digiPH_darkorange, dark]
	\frametitle{Sintesis}

	\textbf{Kesimpulan}\\
	Pemilihan tapak merupakan proses multikriteria yang menghubungkan kondisi lahan dengan kebutuhan perancangan.

	\vspace{0.4cm}
	\textbf{Rangkaian berpikir analitis:}
	\begin{center}
		Lokasi \& luas $\rightarrow$ Bentuk \& topografi $\rightarrow$ Aksesibilitas \\
		$\downarrow$ \\
		Vegetasi \& kondisi tanah $\rightarrow$ Ketampakan \& karakter lingkungan \\
		$\downarrow$ \\
		Penilaian lahan $\rightarrow$ Matriks pemilihan $\rightarrow$ \textbf{Tapak paling sesuai}
	\end{center}
\end{frame}

% ---------------------------------------------------
% FRAME 17: Tugas Analisis (Bagian 2 dari Sintesis & Tugas)
% ---------------------------------------------------
\begin{frame}[mytitle=standard, t, mycolor=digiPH_white, light]
	\frametitle{Tugas Analisis}

	\textbf{Tugas mahasiswa:} Pilih 3 alternatif tapak untuk satu jenis bangunan, kemudian:
	\begin{itemize}
		\item Identifikasi lokasi dan luas.
		\item Analisis bentuk tapak, topografi, aksesibilitas, ketampakan.
		\item Identifikasi vegetasi dan kondisi tanah.
		\item Tentukan kriteria \& bobot, susun matriks pemilihan, hitung skor.
		\item Berikan argumentasi mengapa tapak terpilih paling sesuai.
	\end{itemize}

	\vspace{0.3cm}
	\textbf{Pertanyaan reflektif:}\\
	Apakah tapak dengan skor tertinggi secara numerik selalu menjadi tapak terbaik secara arsitektural? Jelaskan berdasarkan hubungan antara data tapak, tujuan proyek, dan strategi perancangan.

	\vspace{0.2cm}
	\footnotesize{\textit{Basis materi: Kriteria Pemilihan Tapak (termasuk rujukan FIA UK 2018).}}
\end{frame}

\bibliographystyle{apalike}
\bibliography{../filename.bib} % required for citecompletion, not work in mainfile
\end{document}

